\section{Global old and born NC geometry through the permitted finite prephase}
\label{sec:nc-permitted-global-carrier}

Let $M=(X,X,E,I,d)$ have coherent ideal, positive mark and compact
whole controlled cosupport. The original NC carrier $H$ is adapted to
the entire retained SNC list $E$. The conclusions below retain global
NC geometry, including outside the permitted locus, while the asserted
simplicity of the old carrier concerns an actual neighbourhood of the
whole current cosupport. No old selfcrossing outside that neighbourhood
is repaired merely to make the global old carrier simple. This chapter
does not assert old-branch admissibility at every intermediate centre,
identify its phases with the old ordered invariant, or provide the
complete merely-NC initial-boundary resolution and its functor laws.

\begin{theorem}[The original permitted finite prephase preserves the global joint carrier]
\label{hir:nc-original-global-finite-permitted-prephase}
\lean{Hironaka.HironakaNormalCrossingsPermittedGlobalCarrier.exists_original_global_finite_permitted_prephase}
\leanok
The original data supply an actual finite $M$-test tower $T:X_T\to X$
such that $\operatorname{strict}_T H$ is a global NC carrier adapted to
the \emph{entire} current boundary $E_T$. They also supply a literal
open $V\subset X_T$ containing the entire current controlled cosupport.
On this same $V$, the old trace has intrinsic branch count at most one,
is closed and smooth, and is transverse to the entire boundary trace.
Adjoining the old trace to that very list gives an SNC divisor on $V$.
No global SNC assertion for the old carrier is made outside $V$.
\uses{hir:nc-original-phase-whole-cosupport-count-drop,hir:nc-compact-current-cosupport-without-domain,hir:nc-open-image-actual-equivalence,hir:compact-cosupport-open-tower-extension,hir:general-nc-whole-biholomorphic-carrier-pair,hir:general-nc-biholomorphic-intrinsic-multiplicity,hir:nc-boundary-adapted-of-actual-open-traces,hir:nc-boundary-adapted-actual-empty-tower,hir:general-nc-count-one-smooth-terminal-pair}
\end{theorem}
\begin{proof}
First establish the actual global extension argument used in the
construction. Suppose an actual open $U$ contains the whole cosupport
of a well-formed marked object $A$, and a genuine finite test history
$S$ on $U$ has an actual final adapted old carrier and entire boundary.
The actual finite reverse-extension theorem supplies a genuine ambient
test history $Q$ and a whole tower isomorphism from $S$ to $Q|_U$ over
the identity of $U$. Its actual endpoint biholomorphism identifies the
strict old carriers and the entire transformed marked records, hence
every position of the entire boundary lists. Pulling the local adapted
coordinates back along its genuine inverse therefore gives the required
adapted trace on the actual open preimage of $U$.

The complementary part is derived from real testness. Every actual
centre of $Q$ lies over the original whole cosupport of $A$. This
cosupport is closed by coherence and the smooth-carrier order-locus
theorem. The restriction of $Q$ over its actual open complement thus has
only empty centres. The genuine empty-history theorem transports the
original adapted pair to this restriction with every literal empty born
position retained. Exact open-restriction dictionaries identify that
pair with the trace of the ambient strict old carrier and entire ambient
current boundary. The two actual opens cover the whole endpoint: an
original base point outside $U$ lies outside the entire original
cosupport. The ambient strict old carrier and every current boundary
member are closed. The actual open-trace gluing theorem therefore gives
the global adapted pair for $Q$. This argument uses a genuine extended
history, its genuine inverse and its actual empty complementary history;
no desired global tracking conclusion is assumed.

Now prove the asserted construction simultaneously for every ambient
manifold and every global old-count bound $r$, by strong induction on
$r$, retaining both the global adapted pair and the final permitted
neighbourhood. If $r\le1$, use the identity tower and the whole ambient
open. The original adapted pair supplies the global assertion; exact
open transport gives the same intrinsic count on that open.

For $r>1$, the actual original coefficient-stratum phase supplies a
finite original test tower $P$, its global current adapted pair, and a
literal open $U$ containing its entire current original cosupport on
which old count is at most $r-1$. Actual properness and the current
closed order locus make that whole current cosupport compact. Its
literal restriction to $U$ is compact because its image under the open
inclusion is exactly the entire ambient compact cosupport. The current
ideal is coherent, the positive mark is unchanged, and all current
boundary entries are included in the actual restricted object. Apply
the induction hypothesis to these actual restricted data, obtaining a
finite local suffix $S$, its \emph{global} adapted pair on its entire
local endpoint, and its final whole-cosupport neighbourhood.

Apply the genuine reverse-extension producer to this very $S$. The
preceding two-open argument, using the actual global pair after $P$ and
the recursively produced global pair after $S$, proves that its genuine
ambient extension $Q$ has a global adapted current pair. Every current
permitted point after $Q$ projects into the whole current cosupport
after $P$, hence into $U$. The whole marked-record equality places its
unique actual local representative in the local current cosupport and
therefore in the recursively chosen local open. The actual upper open
embedding takes this open to a literal ambient open $W$ containing the
\emph{entire} current cosupport. Its genuine restricted inverse
identifies the old carrier and each boundary entry by the point identity
$g(\chi^{-1}w)=w$. The actual biholomorphic intrinsic count law transfers
the bound by one to this same $W$.

Append the genuine towers $P$ and $Q$. Exact append dictionaries
identify their whole current marked record and literal strict original
carrier, giving both the global adapted pair and the same final open
count bound. Since $r-1<r$, the induction is finite. Actual NC branch
counts are bounded by the ambient dimension, which supplies the initial
$r$ without a caller-supplied bound. Finally, actual count-one
coordinates make the old trace a closed smooth submanifold transverse
to every boundary entry, and give SNC after adjoining it on this very
open. The public construction has only the stated original data as
inputs; all phases, histories, extensions and geometric tracking are
internally constructed.
\end{proof}

\begin{corollary}[Original NC geometry supplies the empty-retained-boundary case]
\label{hir:nc-original-empty-retained-global-finite-prephase}
\lean{Hironaka.HironakaNormalCrossingsPermittedGlobalCarrier.exists_original_nc_global_finite_permitted_prephase}
\leanok
For an original NC carrier $H$, coherent positive-mark ideal and compact
whole cosupport, the same actual finite tower exists with initial
retained boundary empty. Its literal strict old carrier remains globally
adapted to the complete born boundary, with every empty position
retained. On the actual produced open containing the entire current
controlled cosupport, adjoining the old trace gives SNC.
\uses{hir:nc-original-global-finite-permitted-prephase,hir:general-nc-boundary-empty-producer}
\end{corollary}
\begin{proof}
The original NC coordinates themselves give the actual adapted pair
$(H,[])$. Apply the theorem to this very pair and the unchanged original
ideal and mark. The actual transform of the empty retained list is
precisely the entire accumulated born list; no filtering removes any
position. Thus its global adapted-pair conclusion and its same-open
count-one conclusion are exactly the claimed ones. This does not replace
the separate initial labelled-NC boundary and whole-history centre
admissibility obligations by the hypothesis of an already simple old
boundary.
\end{proof}
